\chapter{Methods}
\label{ch:methods}

\section{Configuration System}
Experiments are implemented in PyTorch \parencite{pytorch_2019} using PyTorch Lightning \parencite{pytorch_lightning_2025}.
Configuration is managed through composable YAML files, separating shared defaults (backbone, optimizer, data loading) from track-specific settings (augmentations, loss functions).
An individual run combines a base configuration, a track configuration, and an optional loss configuration, with specific parameters overridden via command line arguments.

\section{Model Architecture}
A 1D adaptation of the ResNet-18 architecture \parencite{resnet_2015} that operates directly on 12-lead ECG waveforms (Figure~\ref{fig:model_arch}) was implemented.
The network consists of an initial convolutional layer followed by four residual stages, each containing two blocks with 1D convolutions and internal temporal downsampling.
Global average pooling aggregates the features into a fixed-size vector.
Depending on the track, this backbone is followed by either a classification head (for classification learning) or a projection head (for contrastive learning).

\begin{figure}[htbp]
  \centering
  \resizebox{0.6\textwidth}{!}{%
    \begin{tikzpicture}[node distance=4mm,>=latex]
      \tikzset{
        arrow/.style={-Latex, thick},
        annot/.style={font=\footnotesize},
        op/.style={draw, rounded corners=2pt, thick, align=center, minimum height=5mm, minimum width=18mm, fill=gray!10},
        stem/.style={draw, rounded corners=2pt, thick, align=center, minimum height=5mm, minimum width=22mm, fill=gray!10},
        blk/.style={draw, rounded corners=2pt, thick, align=center, minimum height=5mm, minimum width=18mm},
        head/.style={draw, rounded corners=2pt, thick, align=center, minimum height=8mm, minimum width=18mm},
        tiny/.style={draw, rounded corners=2pt, thick, align=center, minimum height=4.5mm, minimum width=16mm, font=\footnotesize},
      }
      % =========================================================
      % Overview
      % =========================================================
      \node[op] (x) {12-lead ECG};
      \node[stem,below=3mm of x] (conv) {Conv-BN-ReLU\\MaxPool};
      \node[blk,below=3mm of conv,fill=orange!20] (resblk) {ResBlk $\times 4$};
      \node[op,below=3mm of resblk] (drop) {GAP-Flatten\\Dropout};

      \draw[arrow] (x) -- (conv);
      \draw[arrow] (conv) -- (resblk);
      \draw[arrow] (resblk) -- (drop);

      % --- split to heads without overlap ---
      \node[head,below=7mm of drop, xshift=-12mm, anchor=north,fill=blue!30] (proj) {Projection\\Head};
      \node[head,below=7mm of drop, xshift= 13mm, anchor=north,fill=blue!30] (clf)  {Classification\\Head};
      \draw[arrow] (drop) -- (proj);
      \draw[arrow] (drop) -- (clf);

      % =========================================================
      % Residual block (zoom)
      % =========================================================
      \node[draw, rounded corners=3pt, thick, inner sep=5pt, fill=orange!20,
        right=8mm of resblk, yshift=+3mm,
        minimum width=27mm, minimum height=57mm] (inset) {};
      \node[annot,font=\bfseries,anchor=north] (t_res)
      at ($(inset.north)+(0,-0.2mm)$) {ResBlk};

      \node[tiny,fill=blue!10,below=8mm of inset.north] (c1) {Conv $3{\times}1$};
      \node[tiny,fill=blue!10,below=3mm of c1] (bn1) {BN};
      \node[tiny,fill=red!20,below=3mm of bn1] (r1) {ReLU};
      \node[tiny,fill=blue!10,below=3mm of r1] (c2) {Conv $3{\times}1$};
      \node[tiny,fill=blue!10,below=3mm of c2] (bn2) {BN};
      \node[tiny,fill=red!20,below=3mm of bn2] (rout) {ReLU};

      \coordinate (inflow) at ($(c1.north)+(0,8mm)$);
      \draw[arrow] (inflow) -- (c1.north);
      \draw[arrow] (c1) -- (bn1);
      \draw[arrow] (bn1) -- (r1);
      \draw[arrow] (r1) -- (c2);
      \draw[arrow] (c2) -- (bn2);
      \draw[arrow] (bn2) -- (rout);

      % skip connection
      \draw[very thick,-Latex]
      ($(inflow)+(0,-1.5mm)$) |- ($(inflow)+(-12mm,-1.5mm)$) |- ($(rout.west)+(-1.5mm,0)$) |- (rout.west);
    \end{tikzpicture}
  }%
  \caption[Model architecture]{Model architecture with classification and projection heads (left) and the residual block definition (right).}
  \label{fig:model_arch}
\end{figure}

To limit the experimental scope, the architecture is fixed across all runs; batch normalization is used. This ensures that performance differences can be attributed to preprocessing, augmentations, and objective functions rather than model capacity.

\section{Data Regimes}
\label{sec:data_regimes}
We investigate three preprocessing levels to isolate the effect of harmonization:
\begin{enumerate}
  \item \textbf{bp}: Band-pass filtered signals ($0.67$--$45$\,Hz).
  \item \textbf{bp\_sc}: Band-pass filtering plus soft clipping to handle outliers.
  \item \textbf{bp\_sc\_norm}: Band-pass filtering, soft clipping, and robust per-lead normalization (z-score using median and IQR).
\end{enumerate}
Experiments read preprocessed tensors directly from disk for efficiency.

\paragraph{Band-pass filtering.}
Filtering is implemented as a linear-phase FIR band-pass using \texttt{scipy.signal.firwin}, applied with zero-phase \texttt{filtfilt} to avoid phase distortion.
Cutoffs are $0.67$--$45$\,Hz (motivated by the Zhao et al.\ ECG quality work and the defaults used in NeuroKit2/BioSPPy-style preprocessing) \todo{Add citations for Zhao et al.\ and NeuroKit2 / BioSPPy}.
The tap length is set to $\lfloor 1.5 \cdot f_s \rfloor$ with odd length enforced ($f_s=400$\,Hz after resampling); taps are reduced for short signals to fit the record length.

\paragraph{Soft clipping.}
Soft clipping is a smooth saturation transform \texttt{softclip}(x; a,b,c) applied per lead.
Lower/upper bounds $(a,b)$ are estimated from training folds only (0--3) by computing per-record $p_1/p_{99}$ and then taking the global 5th percentile of $p_1$ and 95th percentile of $p_{99}$ per lead.
The shape parameter is $c=\log(2)/\left(2/(e^2+1)\right)$ (as implemented in code), and each lead is scaled by $\max(|a|,|b|)$ to keep amplitudes comparable across leads.
\todo{find citation for this, Nabil sent me an article}

\paragraph{Robust per-lead normalization.}
The \textbf{bp\_sc\_norm} regime applies per-record, per-lead robust normalization:
\[
  x_i \leftarrow \frac{x_i - \mathrm{median}(x_i)}{\mathrm{IQR}(x_i) + \epsilon}, \quad \epsilon=10^{-6},
\]
where $\mathrm{IQR} = p_{75}-p_{25}$. No population-level statistics are used.

\section{Augmentations}
\subsection{Default augmentations}
Training uses mild augmentations to improve robustness without distorting ECG morphology:
\begin{itemize}
  \item \textbf{Amplitude scaling:} multiplicative gain jitter. \textbf{bp} uses $[0.95,1.05]$, \textbf{bp\_sc} uses $[0.975,1.025]$, and \textbf{bp\_sc\_norm} disables scaling ($[1.0,1.0]$) to preserve the intent of robust normalization.
  \item \textbf{Gaussian noise:} additive noise with $\sigma=0.005$ (per-view by default).
\end{itemize}
Stronger transforms are explicitly disabled in the main experiments: time warping (max\_time\_warp $=0$), time masking (max\_mask\_duration $=0$), channel masking (mask\_prob $=0$), and baseline-wander augmentation (wandering amplitude $=0$). This preserves interpretability and avoids injecting artifacts not supported by the target clinical setting.
\todo{Add citations for augmentation choices in ECG (e.g., mild gain/noise vs. masking/warping).}

\subsection{Physiological augmentation (rotation)}
For the \textbf{bp} regime only, we optionally apply a random rotation of the frontal electrical axis. This augmentation simulates valid physiological variation by projecting the cardiac dipole onto the limb leads at a slightly perturbed angle (implementation: approximate frontal-axis rotation in VCG space and mapping back to 12-lead ECG).
Axis rotation is enabled only in the \texttt{rot10} runs (maximum absolute rotation $10^\circ$, applied with probability 1), and kept off for \textbf{bp\_sc} and \textbf{bp\_sc\_norm} because the nonlinear preprocessing steps can break the linearity assumptions required for a physically meaningful rotation.
\todo{insert a graphic here how the signal changes per lead when rotated}

\section{Experiment Tracks}
\label{sec:track-definitions}
\subsection{Track 1: Supervised Classification}
Track~1 trains a binary classifier end-to-end. We evaluate combinations of loss functions (weighted BCE, Focal Loss, RAT), preprocessing regimes, and augmentation strategies.
Performance is primarily measured using the challenge-aligned TPR at 5\% FPR, alongside AUROC and Average Precision (AP).
We also compare against a random forest baseline trained on simple per-channel mean/std statistics to quantify the value of learning from raw waveforms.

\subsection{Track 2: Representation Learning}
Track~2 focuses on optimizing the geometric structure of the embedding space without direct classification supervision. We compare:
\begin{itemize}
  \item \textbf{SupCon (standard)}: A baseline supervised contrastive loss.
  \item \textbf{SupCon (imbalanced)}: A variant incorporating imbalance-handling logic.
  \item \textbf{Prototype-based}: An objective using learnable class prototypes.
\end{itemize}
Embedding quality is monitored using the TTC metrics defined in Section~\ref{sec:ttc-metrics} (SAD, SAA, CAC, CAD).

\subsection{Track 3: Linear Probing}
Track~3 evaluates the representations learned in Track~2 by freezing the encoder and training a simple linear classifier on top. This isolates the quality of the learned features from the end-to-end optimization capability of Track~1.

\section{Dataset Splits and Evaluation}
Given the lack of a suitable external dataset, we use a stratified 5-fold cross-validation scheme.
Splits are stratified by both patient and dataset source to prevent leakage and ensure consistent label distributions.
Folds 0--3 are used for training and validation, while fold 4 is strictly held out for final testing.

\section{Analysis Pipeline}
All post-hoc analysis is performed on the held-out test set (fold 4).
The evaluation pipeline generates:
\begin{itemize}
  \item \textbf{Clinical Metrics:} Classification performance (AUROC, TPR@5\%, AP) is calculated on the full test fold.
  \item \textbf{Embedding Diagnostics:} A fixed subset of the test fold (the "probe set") is used to compute geometric metrics (CAC, CAD) and visualize the latent space using UMAP and PCA.
  \item \textbf{Consistency Analysis:} We measure the agreement between different models using ranking correlations and intersection-over-union of the top-5\% predictions.
\end{itemize}
This ensures that all reported comparisons between tracks and preprocessing methods are based on identical test data and metric definitions.

\section{ST-DFT-LRP for XAI}
\paragraph{Implementation source.}
ST-DFT-LRP is implemented using the official DFT-LRP repository (integrated as a submodule) to ensure fidelity to the method definition.
\todo{Insert repository citation and URL.}
The method is applied to model relevance in the time domain and mapped into a short-time frequency grid using window width $H$ and shift $D$.
For interpretability and comparability, the analysis uses non-overlapping windows ($D=H$) \todo{Confirm the exact window width used (e.g., 125 samples at 400\,Hz).}

\paragraph{Qualitative per-sample composites.}
For individual ECGs, we visualize a composite view consisting of:
(i) the ST-DFT magnitude with a signed relevance overlay,
(ii) the time-domain signal with signed relevance overlay, and
(iii) the frequency spectrum with signed relevance overlay.
These panels provide a compact visual summary of where relevance concentrates in time and frequency.
\missingfigure{Composite ST-DFT-LRP visualization for representative samples from each dataset.}

\paragraph{Quantitative aggregation into 4$\times$4 grids.}
To enable comparisons across samples and groups, ST-DFT-LRP relevance is aggregated into a fixed 4$\times$4 matrix:
\begin{itemize}
  \item \textbf{Time structure (rows):} P, QRS, T, and Between-beats (assigned via NeuroKit2 delineation with midpoint-based disambiguation to ensure full coverage).
  \item \textbf{Frequency bands (columns):} 0.67--4\,Hz, 4--12\,Hz, 12--25\,Hz, 25--45\,Hz (consistent with preprocessing range).
\end{itemize}
Within each cell, relevance is summarized using a signed-balance statistic
\[
  \mathrm{balance} = \frac{\sum R}{\sum |R| + \epsilon},
\]
yielding values in $[-1,1]$ that encode directional relevance while normalizing magnitude across beats.
Beat-level matrices are normalized and averaged across beats; per-sample matrices are aggregated across leads using a relevance-mass-weighted mean.
\todo{Cite the exact delineation method (e.g., NeuroKit2 DWT) and any QC filtering applied.}
\missingfigure{Example 4$\times$4 relevance grid and group-mean comparisons.}

\paragraph{Lead diagnostics.}
For each sample, lead relevance mass is summarized to compute lead-importance statistics (per-lead mass distribution, entropy, and top-$k$ leads).
These provide a lightweight check for lead dominance effects and are used as descriptive diagnostics rather than diagnostic claims.
